% !TeX program = lualatex
% Copyright (C) 2026 IEEE UFG Computer Society contributors.
%
% This work may be distributed and/or modified under the conditions of the
% LaTeX Project Public License, either version 1.3c of this license or (at your
% option) any later version. The latest version is available at
% https://www.latex-project.org/lppl.txt.
%
% This work has the LPPL maintenance status `maintained'. The Current
% Maintainer is the IEEE UFG Computer Society. The work consists of the files
% listed in MANIFEST.md. See LICENSE for details and artwork exclusions.
\documentclass[section-numbering=two-digit]{ieeeufgcs}
\usepackage[brazil]{babel}

\IEEECSSetup{
  document-type={Plano de trabalho},
  title={Título do plano},
  subtitle={Período ou ciclo de referência},
  running-title={Plano de trabalho},
  institution={Capítulo Estudantil da IEEE Computer Society na UFG},
  unit={Unidade acadêmica ou organização responsável},
  author={Equipe responsável},
  date={\today},
  location={Goiânia},
  identifier={Versão 1.0}
}

\begin{document}
\maketitle

\section{Apresentação}

Apresente a organização, o contexto do plano e sua relação com a comunidade.
Delimite o período, o público e o escopo institucional do documento.

\section{Motivação}

Descreva o problema ou a oportunidade que justifica o plano. Relacione a
proposta com necessidades observáveis e com a missão da organização.

\section{Expectativas}

Registre resultados verificáveis para o primeiro ciclo.

\begin{itemize}
  \item formação e desenvolvimento das pessoas participantes;
  \item entregas técnicas ou institucionais;
  \item continuidade, documentação e transferência de conhecimento.
\end{itemize}

\begin{IEEECSNotice}[Condição de validade]
Nenhuma meta deve prevalecer sobre segurança, integridade, acessibilidade,
capacidade das pessoas voluntárias ou disponibilidade real de recursos.
\end{IEEECSNotice}

\section{Missão}

Declare a finalidade permanente da iniciativa e os princípios que orientam
sua execução.

\section{Objetivos}

\subsection{Objetivo geral}

Defina o resultado amplo que o plano pretende alcançar.

\subsection{Objetivos específicos}

\begin{itemize}
  \item converter a missão em resultados observáveis;
  \item estabelecer práticas documentadas e repetíveis;
  \item desenvolver responsáveis e pessoas sucessoras.
\end{itemize}

\section{Escalabilidade}

Defina critérios de passagem antes de ampliar o público, os recursos ou a
complexidade operacional.

\begin{IEEECSBrandedTable}{
  >{\raggedright\arraybackslash}p{0.20\linewidth}
  >{\raggedright\arraybackslash}X
  >{\raggedright\arraybackslash}X
}
\IEEECSHeaderCell{Etapa} &
\IEEECSHeaderCell{Critério de passagem} &
\IEEECSHeaderCell{Entrega principal} \\
Fundação &
Responsáveis, regras e recursos iniciais confirmados. &
Plano aprovado e calendário inicial. \\
Piloto &
Primeiras atividades testadas e carga da equipe conhecida. &
Piloto documentado e retrospectiva. \\
Consolidação &
Rotina estável, documentação atualizada e sucessão preparada. &
Relatório do ciclo e transição de gestão. \\
\end{IEEECSBrandedTable}

\section{Sucessão}

Explique como responsabilidades, documentos, acessos e conhecimento serão
transferidos. Identifique responsáveis principais e pessoas em formação.

\section{Projetos e engajamento}

Organize as frentes de trabalho e associe cada uma a uma entrega mínima.

\begin{IEEECSBrandedTable}{
  >{\raggedright\arraybackslash}p{0.24\linewidth}
  >{\raggedright\arraybackslash}X
  >{\raggedright\arraybackslash}X
}
\IEEECSHeaderCell{Frente} &
\IEEECSHeaderCell{Propósito} &
\IEEECSHeaderCell{Primeiro ciclo} \\
Formação &
Desenvolver competências técnicas e organizacionais. &
Atividades, materiais e evidências de aprendizagem. \\
Projetos &
Transformar necessidades delimitadas em entregas verificáveis. &
Protótipo ou relatório acompanhado de documentação. \\
Comunidade &
Integrar participantes e manter comunicação recorrente. &
Calendário, acolhimento e evento de referência. \\
\end{IEEECSBrandedTable}

\subsection{Roteiro do ciclo}

\begin{IEEECSBrandedTable}{
  >{\raggedright\arraybackslash}p{0.18\linewidth}
  >{\raggedright\arraybackslash}p{0.25\linewidth}
  >{\raggedright\arraybackslash}X
}
\IEEECSHeaderCell{Período} &
\IEEECSHeaderCell{Prioridade} &
\IEEECSHeaderCell{Entregas} \\
M1--M3 & Fundar e testar & Aprovar o plano e executar o primeiro piloto. \\
M4--M6 & Validar e abrir & Revisar materiais e ampliar a participação. \\
M7--M9 & Integrar & Executar projetos e registrar resultados. \\
M10--M12 & Transferir & Avaliar o ciclo e preparar a sucessão. \\
\end{IEEECSBrandedTable}

\section{Equipe responsável}

\begin{IEEECSBrandedTable}{
  >{\raggedright\arraybackslash}p{0.30\linewidth}
  >{\raggedright\arraybackslash}X
  >{\raggedright\arraybackslash}X
}
\IEEECSHeaderCell{Função} &
\IEEECSHeaderCell{Responsável} &
\IEEECSHeaderCell{Observações} \\
Coordenação & & \\
Secretaria & & \\
Responsabilidade técnica & & \\
Orientação docente & & \\
\end{IEEECSBrandedTable}

\section{Justificativa}

Reúna os alinhamentos técnico, acadêmico, regional e institucional. Demonstre
que a escala proposta é compatível com os recursos e riscos conhecidos.

\section{Pedido de aprovação}

Registre de maneira direta quais decisões são solicitadas.

\begin{enumerate}
  \item aprovar o plano e seu período inicial;
  \item confirmar responsáveis e condições institucionais;
  \item autorizar as primeiras atividades previstas.
\end{enumerate}

\subsection{Primeiros passos após a aprovação}

Liste as ações imediatas, seus responsáveis e os prazos iniciais.

\section{Conclusão}

Retome o valor esperado, as condições de execução e a decisão solicitada.

\end{document}
